%% notes-tex/database/pputida-cell-state/sections/transposon.tex
%% [[experiments.database.pputida-cell-state]]
%% https://github.com/Mjvolk3/torchcell/tree/main/notes-tex/database/pputida-cell-state/sections/transposon.tex
%% SOURCE: hand-authored from the served records and the loader docstrings cited inline.

\section{What a transposon record means here\secstatus{todo}}\label{sec:transposon}

Four rows of Table~\ref{tab:pcsrows} are transposon screens and they hold almost every
record in it, so the representation is worth stating exactly.

\notesec{The experiment} A randomly barcoded transposon library is a pool in which
each cell carries one transposon inserted at a random genomic position, with a random
DNA barcode in the insert. An insertion inside a gene disrupts it, so each cell is a
loss-of-function mutant of one gene. The pool grows competitively in a condition,
barcodes are sequenced at time zero and at the end, and a gene's fitness is the log2
ratio of its barcode abundance after growth over time zero, averaged across that gene's
insertion strains and normalized so the typical gene sits at zero.

\notesec{The genotype is a gene, not a clone} Each record carries one
\sym{TransposonInsertionPerturbation} with \texttt{insertion\_position},
\texttt{insertion\_strand} and \texttt{barcode} left null. Those nulls are the
representation, not missing data: the stored value is the gene-level aggregate over
every insertion mutant of that gene, which is the quantity the release reports. Royet
2025 is the case that makes this unavoidable, since its mariner library is not barcoded
at all and a clone is unresolvable in principle.

\notesec{The phenotype is not the fitness class} These records use
\sym{EnvironmentResponsePhenotype}, whose measurement type is \texttt{log2\_ratio} and
whose assay type is \sym{pooled\_competitive\_growth\_barcode}.
\sym{FitnessPhenotype} would be the obvious class and it is wrong here:
\texttt{validate\_fitness} clamps a non-positive value to 0.0, which would erase every
fitness defect in the compendium, since a defect is exactly a negative log2 ratio. The
signed ratio therefore goes on the environment-response class.

\notesec{Reliability travels with the value} The Fitness Browser row carries a
moderated $t$ in \sym{environment\_response\_t\_statistic}, and the loader derives a
standard error as $\sqrt{(f/t)^2 - 0.1^2}$, where $f$ is the fitness value. That makes
it the one compendium row where a per-record weight is available, which matters for any
model that should down-weight a gene measured in a sparse pool.

\notesec{Why these rows are a feature source and not labels} For the isoprenol task of
Sec.~\ref{sec:isoprenol}, three differences rule them out as training labels.
The readout is competitive growth, not product titer, and a gene that raises titer can
be fitness-neutral or fitness-costly. The perturbation differs in kind: a transposon
insertion is a full disruption while CRISPRi is partial repression of a gene that is
still present, so the sign may agree but the magnitudes share no scale. And the chassis
differs, since the libraries are wild-type KT2440 and carry neither the mevalonate
pathway nor its flux burden.

What they can supply is a descriptor. A gene's fitness across the compendium's many
carbon and nitrogen sources is a cheap, genome-wide measure of how essential or
buffered that gene is, available for nearly every CRISPRi target without a new
experiment. \emph{Hypothesis, untested:} that descriptor improves last-round titer
prediction over the proteome features alone. Nothing here measures it, and it should be
checked against a proteome-only model before it is relied on.
